% Source pins (SHA-256): PifoStatement.lean=fe9475ebfa2c462d27cb0a99781926074e1281396f2c9e8d32d6dfca788d814e; pifo-elegant-v4.md=356ea38146ced437e20200db04e2ce220c73f9ac9fa892b746f783b8a2d49ab8.
\section{Flushes determine every interleaving}
\label{sec:main-proof}

We prove \autoref{thm:main}.  It is convenient to work over an arbitrary finite
set \(A\); transporting schedulers and operation words along a bijection
\(A\cong\Fin(k)\) gives the frozen statement.

For \(B\subseteq A\), write \(\restrict{w}{B}\) for the subsequence of a word
\(w\) whose letters lie in \(B\).  A natural-valued rank function induces a
total preorder.  For such a preorder \(r\), \(\sort_r(w)\) is the stable sort
of the \emph{occurrences} in \(w\) by
\((r(a),\text{source arrival})\).  For a partition \(P\) of \(A\),
\(\proj{P}(w)\) replaces each letter by its \(P\)-block.  If \(M\) is a
scheduler, \(F_M:A^*\to A^*\) denotes its flush transformer, with the
\(\Some\) constructors erased.  Validity ensures that a flush of \(n\) pushes
has \(n\) successful pops.

\subsection{Operational facts}

The following elementary invariant connects the atomic formal semantics with
the queue-level reasoning used later.

\begin{lemma}[Reference balance]
\label{lem:balance}
At every valid reachable state, for every internal node and each of its
children \(c\),
\begin{equation}
 \#\{\text{queued references to }c\}
 =
 \#\{\text{pending packets in the subtree of }c\}.
 \label{eq:balance}
\end{equation}
Consequently a root pop fails exactly when the total pending-packet count is
zero; after an internal queue selects a child, recursive failure is impossible.
\end{lemma}

\begin{proof}
Initially both sides of \autoref{eq:balance} are zero at every node.  A valid
push through child \(c\) appends exactly one reference to \(c\) and exactly one
packet somewhere below \(c\); it does not affect the equality for another
child.  A successful pop selecting \(c\) removes one such reference and,
recursively, one packet below \(c\).  Induction down the selected path gives a
successful leaf pop.  Hence a positive root count supplies a complete
root-to-leaf pop, while count zero means the root queue is empty.  The atomic
partial-failure branch exists in the definition but is unreachable from a
valid reachable state.
\end{proof}


In particular, for a fixed operation word, the \(\None/\Some\) pattern depends
only on the number of preceding pushes and successful pops, not on ranks or
topology.  Two valid schedulers therefore have identical definedness patterns.

\begin{lemma}[Timestamp compression]
\label{lem:compression}
Replace the arrival stamps occurring in any embedded queue computation by
\(1,2,\ldots\) in their existing strict order.  Every queue selection, and
hence every emitted value, is unchanged.
\end{lemma}

\begin{proof}
Ranks are unchanged.  Equal-rank entries are compared only by the order of
their arrival stamps, which a strictly increasing replacement preserves.
Every choice of a least \((\rank,\arr)\) key is therefore the same.  Induction
over the operation word gives the claim.
\end{proof}

This lemma allows a child invoked at sparse global timestamps to be compared
with a standalone run whose pushes are consecutively stamped.  We will also
attach a ghost source-occurrence tag to every entry when couplings need to
identify an insertion.  The tag is not part of the model and is never
observed; erasing it commutes with every transition because queue ordering uses
only rank and arrival.  Within every reachable queue, and within every abstract
pending-occurrence list used below, arrival stamps are pairwise distinct.  The
same source push may deliberately occur at several selector levels with the
same stamp.  Consequently \((\rank,\arr)\) is a strict key within each
selector: every selection has a unique minimum.  Arrival injectivity is used
in the colored-PIFO coupling, while the strict-key consequence is used in the
diamond and collapse arguments.

\subsection{What two types reveal}

\begin{lemma}[Two-type reduction]
\label{lem:two-type}
For distinct types \(x,y\), the restriction of a stateless PIFO tree to
operation words over \(\{x,y\}\) behaves as one stable PIFO with an effective
comparison
\[
  e(x,y)\in\{x<y,\ x\sim_e y,\ y<x\},
\]
where \(\sim_e\) denotes a rank tie rather than equality of types.
The comparison is determined by the two binary flushes in
\autoref{tab:binary-recovery}.
\end{lemma}

\begin{proof}
Follow the assigned paths of \(x\) and \(y\) together.  Stop at their first
internal node with different child indices, or at their common leaf if their
indices never differ.  At every earlier node all relevant references carry the
same child value.  Such a node forwards one anonymous service for each
successful pop and cannot choose between the two types.

At a decisive internal node, its selected child identifies \(x\) or \(y\), and
the selected child's restricted subtree is monochromatic.  At a decisive leaf,
the queue emits \(x\) or \(y\) directly.  In either case one rank comparison at
this node, with FIFO order on equality, determines every restricted operation
word.  If \(x\) is strictly earlier, both two-letter arrival orders drain as
\(xy\); if the ranks tie, each arrival order is retained; and if \(y\) is
strictly earlier, both drain as \(yx\).  These are exactly the three rows of
\autoref{tab:binary-recovery}.
\end{proof}

\begin{table}[H]
\centering
\begin{tabular}{@{}c@{\qquad}cc@{}}
\toprule
effective comparison \(e(x,y)\) & \(F(xy)\) & \(F(yx)\)\\
\midrule
\(x<y\) & \(xy\) & \(xy\)\\
\(x\sim_e y\) & \(xy\) & \(yx\)\\
\(y<x\) & \(yx\) & \(yx\)\\
\bottomrule
\end{tabular}
\caption{Recovery of the effective binary comparison.  Both arrival orders are
needed to distinguish a rank tie from either strict order.}
\label{tab:binary-recovery}
\end{table}

Notice that \(e(x,y)\) is an extensional fact about the restricted scheduler;
it need not be read from a particular syntactic node.  Flush-equivalent
schedulers have the same \(e\) for every pair by
\autoref{tab:binary-recovery}.

\subsection{A colored-PIFO coupling}

The induction will replace a root preorder while retaining only the sequence
of selected children.  The next lemma justifies that replacement for arbitrary
interleavings, including exact cross-color rank ties.

\begin{lemma}[Colored PIFOs]
\label{lem:colored}
Let \(p,q\) be total preorders on a finite set of types, and let
\(\gamma\) color those types.  Suppose
\begin{equation}
 \gamma(x)\ne\gamma(y)
 \quad\Longrightarrow\quad
 \cmp_p(x,y)=\cmp_q(x,y),
 \label{eq:cross-color-agree}
\end{equation}
where comparison has the three values less, equal, and greater.  PIFO queues
ordered by \(p\) and \(q\), subjected to identical pushes and pops, emit
identical color traces.
\end{lemma}

\begin{proof}
Refine each color into \emph{cross-profile classes}.  Define
\begin{equation}
 x\sim y
 \quad\Longleftrightarrow\quad
 \gamma(x)=\gamma(y)
 \ \land\ 
 \forall z\text{ of another color}.\ 
   \cmp_p(x,z)=\cmp_p(y,z).
 \label{eq:profile-class}
\end{equation}
By \autoref{eq:cross-color-agree}, the same classes are obtained from \(q\).
We use three consequences of totality and transitivity.

First, two distinct classes of one color are strictly separated in the same
direction by both preorders.  Indeed, an outside type witnessing the difference
between their profiles lies weakly on different sides of them; transitivity
rules out reversing that separation.  Second, if a class ties a type of another
color, it is \emph{value-pinned}: every member ties that witness and therefore
all members tie one another, under both \(p\) and \(q\).  Third, comparisons
between differently colored occurrences agree outright by hypothesis.

Couple the two queues while maintaining the following invariant:
\begin{enumerate}
\item every profile class has the same pending count on both sides; and
\item every class that ties another class has literally the same pending list
      of ghost-tagged occurrences on both sides.
\end{enumerate}
The invariant holds initially and identical pushes extend both copies in the
same class.

Consider a pop.  Suppose the \(p\)-queue selects occurrence \(o_1\) in class
\(K_1\), while the \(q\)-queue selects \(o_2\) in a different class \(K_2\).
Equal class counts provide a pending \(K_2\)-occurrence on the \(p\)-side and a
pending \(K_1\)-occurrence on the \(q\)-side.  Minimality of the selected
entries implies, member-wise,
\[
  K_1\le_p K_2
  \qquad\text{and}\qquad
  K_2\le_q K_1.
\]
The classes cannot have one color: their strict, same-direction separation
would contradict these inequalities.  They have different colors, so their
constant cross-class comparison is common to \(p\) and \(q\).  The two
inequalities force that comparison to be a tie in both directions.

Both classes are therefore value-pinned.  By the invariant, each class's
pending occurrence list is identical across the queues.  Thus \(o_1\) and
\(o_2\) are pending in both queues, and their rank components tie within each
queue.  Minimality on the \(p\)-side gives
\(\arr(o_1)\le\arr(o_2)\); minimality on the \(q\)-side gives the reverse, so
the stamps are equal.  Both occurrences lie in each queue, where arrival
stamps are injective, hence \(o_1=o_2\), contradicting \(K_1\ne K_2\).

Both pops therefore select one profile class and emit the same color.  Its
pending count decreases on both sides.  If the class ties another class, it is
value-pinned, so both queues remove its arrival-least pending occurrence---the
same occurrence---and the identical-list clause survives.  Other lists are
untouched.  Empty pops agree because total pending counts agree.  Induction on
the common operation word completes the coupling.
\end{proof}

% Source pin (SHA-256): pifo-elegant-v4.md=356ea38146ced437e20200db04e2ce220c73f9ac9fa892b746f783b8a2d49ab8.
\subsection{Proper flush decompositions}
\label{sec:decompositions}

For induction on the alphabet, we expose the first queue that makes a
nontrivial routing decision.

\begin{lemma}[Proper top decomposition]
\label{lem:decomposition}
If \(|A|>1\), every valid stateless scheduler is online-equivalent to a
scheduler of the form
\begin{equation}
 M=(P,r,(M_B)_{B\in P}),
 \label{eq:decomposition-shape}
\end{equation}
where \(P\) is a proper partition of \(A\) into active root children, \(r\) is
the total preorder induced by root ranks, and \(M_B\) is the child scheduler
restricted to types in \(B\).  Its flush transformer satisfies
\begin{align}
 \restrict{F_M(w)}{B}
   &=F_{M_B}(\restrict{w}{B})
    =F_M(\restrict{w}{B}),
    \label{eq:block-projection}\\
 \proj{P}(F_M(w))
   &=\proj{P}(\sort_r(w)).
    \label{eq:root-pattern}
\end{align}
\end{lemma}

\begin{proof}
Starting at the root, follow the unique active child used by all types for as
long as one exists.  Each skipped queue contains only that child value and
therefore forwards exactly one service per successful pop; removing such a
node preserves every online trace.  The walk ends either at the first node with
at least two active children or at a leaf.  In the former case, partition types
by the chosen child.  Every block is nonempty and, because at least two are
active, proper.

In the leaf case, replace the leaf by an internal root having the old leaf rank
for each type and one FIFO singleton child per type.  The new root queue mirrors
the old leaf queue occurrence-for-occurrence, while a singleton child has no
choice to make.  Since \(|A|>1\), the singleton partition is proper.  This
establishes \autoref{eq:decomposition-shape} without changing online behavior.

During a flush, child \(B\) receives all pushes in
\(\restrict{w}{B}\) before it receives service, and the root calls it exactly
\(|\restrict{w}{B}|\) times.  Its complete output is therefore
\(F_{M_B}(\restrict{w}{B})\), proving the first equality of
\autoref{eq:block-projection}.  On a word confined to \(B\), all root references
select \(B\), which proves the second.  Finally the root drains its reference
occurrences by the stable \((r,\arr)\) order.  Replacing their source types by
their child blocks yields \autoref{eq:root-pattern}.
\end{proof}

We call \((P,r)\) a \emph{decomposition} of a word transformer \(F\) when
\autoref{eq:block-projection} and \autoref{eq:root-pattern} hold.  Two consequences will be used
without further comment:
\begin{align}
  F(v)&=F_{M_B}(v) &&\text{if every letter of \(v\) lies in \(B\)},
  \label{eq:block-autonomy-simple}\\
  \restrict{F(w)}{B}&=F(\restrict{w}{B}) &&\text{for arbitrary \(w\)}.
  \label{eq:block-filter-simple}
\end{align}
The first is \emph{block autonomy}; the second says that deleting other blocks
commutes with the transformer.

\subsection{Combining two decompositions}
\label{sec:meet}

Suppose \((P,r)\) and \((Q,s)\) decompose one common flush transformer \(F\).
Their meet partition is the set of nonempty incidence cells
\begin{equation}
 R=P\wedge Q
  =\{B\cap C\ne\varnothing:B\in P,\ C\in Q\}.
 \label{eq:meet-partition}
\end{equation}
We will construct a total preorder that agrees with \(r\) across \(P\)-blocks
and with \(s\) across \(Q\)-blocks.  We do \emph{not} need to claim that the
resulting pair \((R,t)\) itself decomposes \(F\).

First, every cell inherits autonomy.  For \(D=B\cap C\), repeated use of
\autoref{eq:block-filter-simple} gives
\begin{equation}
\begin{aligned}
 \restrict{F(w)}{D}
 &=\restrict{(\restrict{F(w)}{B})}{C}
  =\restrict{F(\restrict{w}{B})}{C}\\
 &=F(\restrict{(\restrict{w}{B})}{C})
  =F(\restrict{w}{D}).
\end{aligned}
\label{eq:cell-autonomy}
\end{equation}

For types in distinct \(R\)-cells, prescribe the effective comparison
\(e(x,y)\) recovered by \lemref{lem:two-type}.  If \(P\) separates the pair, its
two-type restriction is chosen at the \(P\)-root, and therefore
\begin{equation}
 e(x,y)=\cmp_r(x,y).
 \label{eq:e-from-r}
\end{equation}
Likewise, if \(Q\) separates it,
\begin{equation}
 e(x,y)=\cmp_s(x,y).
 \label{eq:e-from-s}
\end{equation}
The prescriptions agree when both partitions separate the pair because both
are recovered from the same two flushes of \(F\).

\begin{figure}[t]
\centering
\begin{tikzpicture}[
  cell/.style={minimum width=2.0cm,minimum height=1.2cm,draw},
  pt/.style={circle,fill=pifogray,text=white,inner sep=2.1pt,font=\small},
  every node/.style={font=\small}]
  \draw[step=2cm] (0,0) grid (4,2.4);
  \node[pt] (x) at (1,1.8) {$x$};
  \node[pt] (y) at (3,1.8) {$y$};
  \node[pt] (z) at (1,.6) {$z$};
  \node[above] at (1,2.4) {$C_1$};
  \node[above] at (3,2.4) {$C_2$};
  \node[left] at (0,1.8) {$B_1$};
  \node[left] at (0,.6) {$B_2$};
  \node[draw=pifored,very thick,rounded corners,fit=(x)(y),inner sep=10pt] {};
  \node[draw=pifoblue,very thick,rounded corners,fit=(x)(z),inner sep=10pt] {};
  \node[pifored,anchor=west] at (4.18,1.78) {$P$-block};
  \node[pifoblue,anchor=east] at (-.18,.32) {$Q$-block};
  \node[align=left,anchor=west] at (5.35,.62)
    {rows expose $\{x,y\}\mid\{z\}$;\\
     columns expose $\{x,z\}\mid\{y\}$};
\end{tikzpicture}
\caption{The only nontrivial three-cell incidence shape.  Three meet cells
form an L when neither all \(P\)-coordinates nor all \(Q\)-coordinates are
distinct.  Its two visible decompositions form the diamond used below.}
\label{fig:meet-diamond}
\end{figure}

\begin{lemma}[Three-cell diamond]
\label{lem:diamond}
On one representative type from each of any three distinct meet cells, the
pairwise comparisons \(e\) form a total preorder.
\end{lemma}

\begin{proof}
View the cells as occupied positions in the \(P\times Q\) incidence grid.  If
the three \(P\)-coordinates are distinct, all comparisons come from the genuine
preorder \(r\) by \autoref{eq:e-from-r}.  If their \(Q\)-coordinates are distinct,
they all come from \(s\) by \autoref{eq:e-from-s}.  Otherwise the three positions
form an L as in \autoref{fig:meet-diamond}; naming its corner \(x\), the restricted
transformer has both decompositions
\begin{equation}
  \{x,y\}\mid\{z\},
  \qquad
  \{x,z\}\mid\{y\}.
  \label{eq:l-decompositions}
\end{equation}

Assume that the non-strict relation induced by \(e\) is not transitive.  Since
each pair is totally compared, there are distinct \(a,b,c\) with
\begin{equation}
  a\le_e b,
  \qquad b\le_e c,
  \qquad c<_e a.
  \label{eq:bad-triple}
\end{equation}
Push one occurrence of each in the cycle order \(a,b,c\), then flush.  The
arrival stamps are distinct, so every queue computation has a unique least
\((\rank,\arr)\) key.  The cycle inequalities determine every strict or weak
root comparison, and the fixed arrival order resolves any remaining rank tie;
no separate tie case split is needed.

Consider in full the L orientation whose grouped pairs are \(\{b,c\}\) and
\(\{c,a\}\).  In the decomposition
\(\{b,c\}\mid\{a\}\), separated-pair comparisons give root ranks
\(\rho_a\le\rho_b\) and \(\rho_c<\rho_a\).  Thus
\(\rho_c<\rho_b\), and the root drains references in the order
\(c,a,b\): if \(\rho_a=\rho_b\), arrival puts \(a\) first.  The resulting
child slots are
\[
  \{b,c\},\ a,\ \{b,c\}.
\]
Within the pair child, \(b\) arrived before \(c\) and \(b\le_e c\), so the
binary recovery table says that its output is \(b,c\).  The complete output is
therefore \(b,a,c\).

In the decomposition \(\{c,a\}\mid\{b\}\), separated comparisons give
\(\sigma_a\le\sigma_b\le\sigma_c\).  With arrival order \(a,b,c\), the root
drains references in that same order, producing slots
\[
  \{c,a\},\ b,\ \{c,a\}.
\]
The pair child outputs \(c,a\), since \(c<_e a\), and the complete output is
\(c,b,a\).  The first letters differ, contradicting that both decompositions
compute the same \(F(abc)\).

Cyclic renaming covers the other two placements of the two grouped edges in
the comparison cycle.  Exchanging the roles of the two decompositions covers
the three reflected L orientations.  These exhaust the incidence shape, so no
failure of transitivity is possible.  Reflexivity and totality are already
built into the pair comparisons; hence the restriction is a total preorder.
\end{proof}

\begin{figure}[t]
\centering
\begin{tikzpicture}[
  panel/.style={draw,rounded corners,minimum width=6.1cm,minimum height=2.85cm,
                align=left,inner sep=8pt},
  title/.style={font=\bfseries}]
  \node[panel] (l) at (-3.35,0) {};
  \node[title,anchor=north] at (l.north) {$\{b,c\}\mid\{a\}$};
  \node[align=left,anchor=north west] at ($(l.north west)+(.18,-.48)$)
    {$\rho_c<\rho_a\le\rho_b$\\
     root sources: $c,a,b$\\
     slots: $[\{b,c\}], [a], [\{b,c\}]$\\
     inner output: $b,c$\\
     final output: $\mathbf{b,a,c}$};
  \node[panel] (r) at (3.35,0) {};
  \node[title,anchor=north] at (r.north) {$\{c,a\}\mid\{b\}$};
  \node[align=left,anchor=north west] at ($(r.north west)+(.18,-.48)$)
    {$\sigma_a\le\sigma_b\le\sigma_c$\\
     root sources: $a,b,c$\\
     slots: $[\{c,a\}], [b], [\{c,a\}]$\\
     inner output: $c,a$\\
     final output: $\mathbf{c,b,a}$};
  \draw[-{Stealth},very thick,pifored] ($(l.south)+(0,-.2)$) --
    node[below,font=\small] {contradict a common flush transformer}
    ($(r.south)+(0,-.2)$);
\end{tikzpicture}
\caption{The representative diamond computation for
\(a\le_e b\le_e c<_e a\), using one common arrival order \(a,b,c\).
Internal references choose child slots; the selected pair child may emit a
different source occurrence.}
\label{fig:diamond-computation}
\end{figure}


The diamond controls triples of cells.  The next argument shows that no
higher-arity obstruction remains.

\begin{lemma}[Meet-compatible preorder]
\label{lem:meet-preorder}
If \((P,r)\) and \((Q,s)\) decompose \(F\), there is a natural-valued total
preorder \(t\) such that
\begin{equation}
 \begin{aligned}
  \cmp_t(x,y)&=\cmp_r(x,y)
    &&\text{whenever \(x,y\) lie in distinct \(P\)-blocks},\\
  \cmp_t(x,y)&=\cmp_s(x,y)
    &&\text{whenever \(x,y\) lie in distinct \(Q\)-blocks}.
 \end{aligned}
 \label{eq:meet-compatibility}
\end{equation}
\end{lemma}

\begin{proof}
Build a directed constraint graph on the types.  Between distinct meet cells,
each relation \(x\le_e y\) contributes a weak edge \(x\to y\), marked exactly
when \(x<_e y\).  Thus a preorder tie contributes unmarked edges in both
directions.  It is enough to prove that no directed cycle contains a marked
edge.

Suppose otherwise.  Choose a marked edge \(a\to b\) as the head of such a
cycle and argue by strong induction on its length.  Let \(c,d\) be the next
vertices after \(b\) when present.  Only the three fixed chord positions
\((a,c)\), \((b,d)\), and \((a,d)\) are needed; they are shown in
\autoref{fig:fixed-chords}.

For length two, the return edge \(b\to a\) contradicts the strict comparison
on the marked head.  For length three, the consecutive edges join three
distinct cells, and \lemref{lem:diamond} supplies a genuine total preorder on
the triple; such a preorder cannot contain a marked directed triangle.

Consider a length-four cycle \((a,b,c,d)\).  If \(a\) and \(c\) occupy
distinct cells, compare that chord.  An edge \(c\to a\) closes the marked
triangle \(a\to b\to c\to a\).  Otherwise totality gives \(a\to c\), and it
is marked: in the three-cell preorder, \(a<_e b\le_e c\) implies
\(a<_e c\).  The cycle \(a\to c\to d\to a\) is then a marked triangle.
Likewise, if \(b,d\) occupy distinct cells, an edge \(b\to d\) yields the
marked triangle \(a\to b\to d\to a\), which retains the original marked
head; if that edge is unavailable, totality makes \(d\to b\) strict, yielding
the marked triangle \(b\to c\to d\to b\).  If both diagonals are intra-cell,
the cycle alternates between two meet cells.
All four cross-cell comparisons then come from one genuine preorder: from
\(r\) when \(P\) separates those cells, and otherwise from \(s\), since \(Q\)
must separate them.  Its ranks would have to satisfy
\[
 \rho_a\le\rho_b\le\rho_c\le\rho_d\le\rho_a
 \quad\text{while}\quad
 \rho_b\le\rho_a\text{ is false},
\]
which is impossible.

Now let the length be at least five.  If \(a,c\) lie in distinct cells, an
edge \(c\to a\) closes the marked triangle through \(a,b,c\).  Otherwise the
forward chord \(a\to c\) is strict by the same three-cell argument, and
replacing \(a\to b\to c\) by this marked head shortens the cycle.  Strong
induction excludes it.

We may therefore assume \(a,c\) are in one cell.  If \(b,d\) are in distinct
cells and the forward edge \(b\to d\) exists, replacing
\(b\to c\to d\) by it shortens the cycle while retaining the original marked
head \(a\to b\).  If the forward edge is unavailable, totality makes the
reverse edge \(d\to b\) strict, closing a marked triangle through \(b,c,d\).
Thus \(b,d\) must also be in one cell.

Now \(a,d\) lie in distinct cells: otherwise \(a,c,d\) would share a cell,
contradicting the cross-cell edge \(c\to d\).  If the chord \(d\to a\) exists,
it closes the already excluded length-four marked cycle
\(a\to b\to c\to d\to a\).  In the other orientation, all four cross-cell
comparisons come from one original root preorder because \(a,c\) share one
cell and \(b,d\) share another.  In that preorder,
\(a<_e b\le_e c\le_e d\), so \(a\to d\) is strict.  Replacing the three-edge
prefix by this marked head produces a shorter marked cycle.  Strong induction
excludes every case, so no marked cycle exists.

It remains to realize all constraints by natural ranks.  Write \(z\leadsto x\)
for reflexive graph reachability, allowing a path of length zero, and set
\begin{equation}
 t(x)=
 \#\{z:z\leadsto x\ \land\ x\not\leadsto z\}.
 \label{eq:counting-linearization}
\end{equation}
Let \(L_x=\{z:z\leadsto x\land x\not\leadsto z\}\).  If \(x\to y\) is a weak
edge and \(z\in L_x\), then \(z\leadsto x\to y\).  A return path
\(y\leadsto z\) would imply \(x\leadsto z\), a contradiction; hence
\(z\in L_y\), and \(L_x\subseteq L_y\).  If \(x\to y\) is marked, no path
returns from \(y\) to \(x\), since it would complete a marked cycle.  Thus
\(x\in L_y\), whereas reflexivity gives \(x\notin L_x\).  The inclusion is
strict and finiteness gives \(t(x)<t(y)\).  A prescribed tie contributes weak
edges both ways, so the two inclusions give equal predecessor sets and equal
ranks.  Thus \(t\) realizes every prescribed \(e\)-comparison.  The constraints in
\autoref{eq:e-from-r} and \autoref{eq:e-from-s} give exactly
\autoref{eq:meet-compatibility}.
\end{proof}

\begin{figure}[t]
\centering
\begin{tikzpicture}[
  v/.style={circle,draw,minimum size=6.5mm,inner sep=0pt,font=\small},
  arr/.style={-{Stealth[length=2mm]}},
  chord/.style={dashed,pifoblue,thick}]
  \node[v] (a) at (150:2.15) {$a$};
  \node[v] (b) at (90:2.15) {$b$};
  \node[v] (c) at (30:2.15) {$c$};
  \node[v] (d) at (-30:2.15) {$d$};
  \node[v] (e) at (-90:2.15) {$\cdots$};
  \node[v] (f) at (-150:2.15) {};
  \draw[arr,very thick,pifored] (a) -- node[above left,font=\small,
    text=pifored] {marked head} (b);
  \draw[arr] (b) -- (c);
  \draw[arr] (c) -- (d);
  \draw[arr] (d) -- (e);
  \draw[arr] (e) -- (f);
  \draw[arr] (f) -- (a);
  \draw[chord] (a) -- node[below,font=\scriptsize] {$(a,c)$} (c);
  \draw[chord] (b) -- node[right,font=\scriptsize] {$(b,d)$} (d);
  \draw[chord] (a) to[bend right=18]
    node[below,font=\scriptsize] {$(a,d)$} (d);
  \node[draw,rounded corners,align=left,font=\small,anchor=west]
    at (3.05,.35) {length $4$, both diagonals intra-cell:\
      $a,c\in D_0$ and $b,d\in D_1$\\
      $\Rightarrow$ one genuine preorder governs all edges};
\end{tikzpicture}
\caption{The fixed-head-chord induction.  A longer marked cycle is shortened
or reduced to a three-cell diamond using only the displayed chord positions;
the alternating length-four case is controlled by one original root preorder.}
\label{fig:fixed-chords}
\end{figure}

\subsection{Strong induction: swap, refine, and collapse}
\label{sec:induction}

We now prove \autoref{thm:main} simultaneously for every finite alphabet by
strong induction on \(|A|\).  If \(A\) is empty, every operation is a failed
pop.  If \(|A|=1\), each successful pop emits the sole type, and
\lemref{lem:balance} says that success is determined only by the push/pop count.

Assume \(|A|>1\), and let \(S,T\) be flush-equivalent valid stateless
schedulers.  By \lemref{lem:decomposition}, normalize them online to proper top
decompositions
\begin{equation}
  S=(P,r,(S_B)_{B\in P}),
  \qquad
  T=(Q,s,(T_C)_{C\in Q}),
  \label{eq:two-top-decompositions}
\end{equation}
with common flush transformer \(F\).  Form \(R=P\wedge Q\) and choose the
meet-compatible preorder \(t\) from \lemref{lem:meet-preorder}.

\paragraph{Swap before refining.}
Replace \(S\)'s root order \(r\) by \(t\), leaving its partition and children
unchanged.  Across distinct \(P\)-blocks the two root preorders agree by
\autoref{eq:meet-compatibility}.  Apply \lemref{lem:colored} with colors \(P\).
The root has the same complete online \(P\)-token trace, so every child receives
the same driven word.  Thus
\begin{equation}
  S\intereq S^t
  \qquad\text{and}\qquad
  F_{S^t}=F.
  \label{eq:s-root-swap}
\end{equation}
Symmetrically \(T\intereq T^t\).  Applying
\autoref{eq:root-pattern} directly to the two swapped schedulers gives both pattern
facts, for every \(w\),
\begin{equation}
 \proj{P}(F(w))=\proj{P}(\sort_t(w)),
 \qquad
 \proj{Q}(F(w))=\proj{Q}(\sort_t(w)).
 \label{eq:both-patterns}
\end{equation}
This order of operations is important: after the swaps, no separate meet-cell
control-word construction is needed.

\paragraph{Refine one proper block.}
For every meet cell \(D\in R\), choose a valid scheduler \(U_D\) realizing the
restricted transformer \(F\) on \(D\).  For example, restrict \(S\)'s fixed
assignment to types in \(D\); cell autonomy
\autoref{eq:cell-autonomy} identifies its flush transformer with
\(\restrict{F}{D}\).

Fix \(B\in P\).  Let \(V_B\) have a \(t\)-ordered root routing to the nonempty
cells \(D=B\cap C\) below \(B\), with child \(U_D\) at cell \(D\).  We claim
\begin{equation}
  F_{S_B}=F_{V_B}.
  \label{eq:block-flush-equality}
\end{equation}
Take a word \(v\) over \(B\).  A word over \(B\) is uniquely reconstructed
from its \(Q\)-block pattern and its subsequence in every \(Q\)-block.  By block
autonomy, \(F_{S_B}(v)=F(v)\), whose \(Q\)-pattern is
\(\proj{Q}(\sort_t(v))\) by \autoref{eq:both-patterns}.  The root-pattern equation
for \(V_B\) gives the same pattern.

Fix a block \(C\in Q\) and set \(D=B\cap C\).  If \(D=\varnothing\), both
\(C\)-subsequences are empty because both output words lie in \(B\).  If
\(D\ne\varnothing\), then \(D\in R\), so \(U_D\) is defined.  Because \(v\)
is a \(B\)-word, \(\restrict{v}{C}=\restrict{v}{D}\), and
\begin{equation}
\begin{aligned}
 \restrict{F(v)}{C}
 &=F(\restrict{v}{C})
  =F(\restrict{v}{D})\\
 &=F_{U_D}(\restrict{v}{D})
  =\restrict{F_{V_B}(v)}{C}.
\end{aligned}
\label{eq:block-subsequences}
\end{equation}
Here the first equality is autonomy for the \(Q\)-decomposition, the third is
the choice of \(U_D\), and the last is
\autoref{eq:block-projection} for \(V_B\); filtering by \(C\) and by \(D\) agrees
because every output of \(V_B\) lies in \(B\).  Equal pattern and equal
subsequences prove \autoref{eq:block-flush-equality}.

The partition \(P\) is proper, so \(|B|<|A|\).  The induction hypothesis
applied to \autoref{eq:block-flush-equality} gives
\begin{equation}
  S_B\intereq V_B.
  \label{eq:block-inter-equality}
\end{equation}
Substitute \(V_B\) for every child \(S_B\) under the already swapped
\(t\)-root.  The root reference queue is unchanged.  Child \(B\) receives its
projected pushes and one anonymous pop whenever a \(B\)-token is selected, so
the driven operation word is the same on both sides.  Its global arrival stamps
may have gaps; \lemref{lem:compression} converts them monotonically to the
standalone instance of \autoref{eq:block-inter-equality}.  Child substitution
therefore preserves the entire online trace.

\paragraph{Collapse equal selector levels.}
After substitution, the scheduler has two routing levels: an outer root over
the \(P\)-blocks and, inside each \(B\), a root over its \(R\)-cells.  Both
levels use preorder \(t\).  Let \(U\) instead have one \(t\)-root routing
directly to the same children \(U_D\).

Couple these schedulers with a set \(\Omega\) of pending ghost-tagged
selector-reference occurrences.  Maintain:
\begin{enumerate}
\item the two-level outer queue is \(\Omega\), labeled by \(P\)-block;
\item \(U\)'s root queue is \(\Omega\), labeled by \(R\)-cell;
\item the inner queue below \(B\) is \(\restrict{\Omega}{B}\), labeled by
      \(R\)-cell; and
\item corresponding \(U_D\) states are identical.
\end{enumerate}
Pushes extend exactly the corresponding queues and child state.  On a pop, let
the outer queue select occurrence \(o\in B\).  Every selector uses the same key
\((t,\arr)\), and arrival stamps within \(\Omega\) are pairwise distinct, so
\(o\) is the unique minimum of \(\Omega\) and therefore also of
\(\restrict{\Omega}{B}\).  The inner queue selects the same occurrence; the
one-level root, holding the same \(\Omega\) with the same keys, also selects
\(o\).  There is no full-key tie case.  Both sides call the same \(U_D\) state.
A successful lower pop removes \(o\) from all corresponding selectors.  An
atomic lower failure rolls both sides back identically.  The invariant
survives, and
\begin{equation}
  S\intereq U.
  \label{eq:s-to-u}
\end{equation}

Starting from \(T\), swap its root to \(t\), refine each proper
\(C\in Q\) into a \(t\)-root over the same meet-cell schedulers \(U_D\), and
collapse.  The proof of block flush equality is the mirror image of
\autoref{eq:block-flush-equality}: use the \(P\)-pattern in
\autoref{eq:both-patterns} and autonomy from the \(S\)-side.  It gives
\begin{equation}
  T\intereq U.
  \label{eq:t-to-u}
\end{equation}
Transitivity of equality of run outputs yields \(S\intereq T\).  The argument
has coupled successful outputs; \lemref{lem:balance} restores the identical
\(\None\) positions for over-pops.  This proves \autoref{thm:main}.

\begin{figure}[t]
\centering
\resizebox{.98\textwidth}{!}{%
\begin{tikzpicture}[
  box/.style={draw,rounded corners,align=center,minimum width=2.8cm,
              minimum height=1.0cm,font=\small},
  arr/.style={-{Stealth[length=2mm]},thick},
  lab/.style={font=\scriptsize,align=center,fill=white,inner sep=1.5pt}]
  \node[box] (s) at (0,1.7) {$S$\\$r\to P\to S_B$};
  \node[box] (st) at (4.35,1.7) {$S^t$\\$t\to P\to S_B$};
  \node[box,minimum width=3.6cm] (sr) at (9.0,1.7)
    {$\widehat S$\\$t\to P\to t\to R\to U_D$};
  \node[box] (t0) at (0,-1.7) {$T$\\$s\to Q\to T_C$};
  \node[box] (tt) at (4.35,-1.7) {$T^t$\\$t\to Q\to T_C$};
  \node[box,minimum width=3.6cm] (tr) at (9.0,-1.7)
    {$\widehat T$\\$t\to Q\to t\to R\to U_D$};
  \node[box,minimum width=3.1cm,very thick,pifogreen] (u) at (14.0,0)
    {$U$\\$t\to R\to U_D$};
  \draw[arr] (s) -- node[lab,above=2pt] {colored-PIFO\\swap} (st);
  \draw[arr] (st) -- node[lab,above=2pt] {induction in\\each proper $B$} (sr);
  \draw[arr,pifogreen] (sr) -- node[lab,above right] {strict-key\\collapse} (u);
  \draw[arr] (t0) -- node[lab,below=2pt] {colored-PIFO\\swap} (tt);
  \draw[arr] (tt) -- node[lab,below=2pt] {induction in\\each proper $C$} (tr);
  \draw[arr,pifogreen] (tr) -- node[lab,below right] {strict-key\\collapse} (u);
\end{tikzpicture}%
}
\caption{Proof architecture.  The common preorder \(t\) is installed before
the children are refined.  Strong induction produces the same meet-cell
children on both sides, and identical selector levels collapse to one common
scheduler \(U\).}
\label{fig:proof-architecture}
\end{figure}

\paragraph{Why the induction is well founded.}
Only the original blocks \(B\in P\) and \(C\in Q\) are recursive arguments,
and both top partitions are proper.  A meet cell may equal an entire original
block, and \(V_B\) may consequently have just one active cell, but the measure
is still \(|B|<|A|\), not tree height or the size of an introduced child.
Repeated packet types cause no problem: every projection, stable sort, and
reconstruction above is a computation on occurrences.
